\documentclass{book}
\usepackage[utf8]{inputenc}
\usepackage[T1]{fontenc}
\usepackage{geometry}
\usepackage{graphicx}
\usepackage{float}
\usepackage{hyperref}
\usepackage[english,indonesian]{babel}

\title{Optimasi Pemrosesan Template LaTeX Menggunakan Pendekatan Berbasis AI}
\author{Academic Support Agent}
\date{\today}

\begin{document}
\maketitle

\begin{abstract}

\end{abstract}

\chapter{PENDAHULUAN}


\chapter{TINJAUAN PUSTAKA}


% Original content (commented out):
% % 
%       قالب زیر توسط سید صالح اعتمادی جهت جزوه برداری توسط دانشجویان آماده شده است.
% 		استفاده از این قالب برای عموم با حفظ این دو خط در ابتدای فایل بلا مانع است.
%
% تشکر ویژه از آقای علی حیدری که در تهیه این قالب از قالب‌های لاتک تهیه شده توسط ایشان کمک گرفته شده

\documentclass[oneside]{book}

\def \CourseName {نام\_درس\_جایگزین\_شود}
%\def \CourseName {طراحی و تحلیل الگوریتم‌ها}
%\def \CourseName {برنامه‌سازی پیشرفته}
\def \Instructor {سید صالح اعتمادی}
\def \Semester {نیم‌سال دوم\\سال تحصیلی ۹۸-۹۹}

\input{packages}
\input{commands}
\input{firstpage}


\begin{document}
\pagenumbering{gobble}
\pagestyle{headings}
\maketitle
\pagenumbering{arabic}

\input{session.tex}

\begin{latin}
\bibliography{references}  	
\bibliographystyle{ieeetr}
\end{latin}

\end{document}


\end{document}